% TEMPLATE-MILC-v3.tex - plantilla maestra para todas las guias MILC v3
% Ver scripts/generadores/build-guia-milc.py para la lista completa de
% claves esperadas. El script valida que no queden placeholders sin
% reemplazar antes de escribir el archivo final.

\documentclass[11pt,letterpaper]{article}
\usepackage[margin=1.75cm]{geometry}
\usepackage[table]{xcolor}
\usepackage{fontspec}
\usepackage[spanish]{babel}
\usepackage{tikz}
\usepackage[most]{tcolorbox}
\usepackage{tabularx}
\usepackage{array}
\usepackage{fancyhdr}
\usepackage{titlesec}
\usepackage{ragged2e}
\usepackage{enumitem}
\usepackage{microtype}

\setmainfont{Helvetica Neue}
\setsansfont{Helvetica Neue}
\setlength{\parindent}{0pt}
\setlength{\parskip}{6pt}
\hyphenpenalty=200
\exhyphenpenalty=200
\pretolerance=400
\tolerance=2000
\setlength{\emergencystretch}{1em}
\renewcommand{\arraystretch}{1.25}
\setlist{leftmargin=5mm,itemsep=1mm,topsep=1mm}
\newcolumntype{Y}{>{\RaggedRight\arraybackslash}X}

\definecolor{milcMagenta}{HTML}{E6007E}
\definecolor{milcVino}{HTML}{5A0038}
\definecolor{milcTurquesa}{HTML}{0093A5}
\definecolor{milcVerde}{HTML}{58A923}
\definecolor{milcMostaza}{HTML}{E5B400}
\definecolor{milcOcre}{HTML}{B86600}
\definecolor{milcNegro}{HTML}{241816}
\definecolor{milcGris}{HTML}{F7F7F4}
\definecolor{milcLinea}{HTML}{E8E2DD}
\definecolor{milcGradePrimary}{HTML}{B63E16}
\definecolor{milcGradeDark}{HTML}{5A1F0B}
\definecolor{milcGradeSoft}{HTML}{F8DDD1}
\definecolor{milcGradeAccent}{HTML}{CFE7FF}
\definecolor{milcGradeText}{HTML}{FFFFFF}
\color{milcNegro}

\pagestyle{fancy}
\fancyhf{}
\lhead{\small Tecnología e Informática · Grado 8}
\rhead{\small Guía 27 · MILC v3}
\cfoot{\small\thepage}
\renewcommand{\headrulewidth}{0pt}

\newtcolorbox{titlebox}[1]{
  enhanced,colback=#1,colframe=#1,arc=7mm,boxrule=0pt,
  left=7mm,right=7mm,top=3mm,bottom=3mm,
  before skip=8pt,after skip=8pt,
  fontupper=\sffamily\bfseries\Large\color{white}
}

\newtcolorbox{softbox}[3]{
  enhanced,colback=#2,colframe=#1,borderline west={4pt}{0pt}{#1},
  arc=2mm,boxrule=.4pt,left=4mm,right=4mm,top=3mm,bottom=3mm,
  before skip=6pt,after skip=8pt,title={#3},coltitle=white,
  colbacktitle=#1,fonttitle=\sffamily\bfseries,fontupper=\RaggedRight,
  attach boxed title to top left={xshift=3mm,yshift=-2mm},
  boxed title style={arc=2mm, boxrule=0pt}
}

\newtcolorbox{drawbox}[2]{
  enhanced,colback=white,colframe=#1,arc=4mm,boxrule=1pt,
  left=5mm,right=5mm,top=4mm,bottom=4mm,before skip=8pt,after skip=8pt,
  title={#2},coltitle=white,colbacktitle=#1,fonttitle=\sffamily\bfseries,
  fontupper=\RaggedRight,
  attach boxed title to top left={xshift=4mm,yshift=-2mm},
  boxed title style={arc=3mm, boxrule=0pt}
}

\newtcolorbox{infoband}[1]{
  enhanced,colback=#1!8,colframe=#1!50,arc=2mm,boxrule=.5pt,
  left=4mm,right=4mm,top=2mm,bottom=2mm,before skip=4pt,after skip=4pt,
  fontupper=\small\RaggedRight
}

\newcommand{\linead}[1]{\noindent\color{milcLinea}\rule{#1}{0.5pt}\color{milcNegro}}
\newcommand{\respuesta}[1]{\par\vspace{2mm}\foreach \n in {1,...,#1}{\noindent\color{milcLinea}\rule{\linewidth}{0.5pt}\par\vspace{5mm}}\color{milcNegro}}
\newcommand{\checkbox}{\textcolor{milcGradeDark}{\fbox{\phantom{x}}}\hspace{5pt}}

\newcommand{\phasecard}[4]{
\begin{tcolorbox}[enhanced,colback=#3,colframe=#2,arc=3mm,boxrule=.5pt,height=3.05cm,valign=center,halign=center,left=1.5mm,right=1.5mm,top=2mm,bottom=2mm]
{\sffamily\bfseries\fontsize{22}{24}\selectfont\textcolor{#2}{#1}}\par
{\sffamily\bfseries\small #4}
\end{tcolorbox}
}

\begin{document}
\thispagestyle{empty}
\begin{tikzpicture}[remember picture,overlay]
  \fill[white] (current page.south west) rectangle (current page.north east);
  \path[fill=milcGradePrimary,rounded corners=16mm]
    ([xshift=.55cm,yshift=-.55cm]current page.north west) rectangle
    ([xshift=-.55cm,yshift=3.65cm]current page.south east);

  \node[anchor=north west,text=milcGradeText] at ([xshift=2.0cm,yshift=-1.85cm]current page.north west)
    {\sffamily\bfseries\fontsize{14}{16}\selectfont GRADO};
  \node[anchor=north west,text=milcGradeText,opacity=.82] at ([xshift=2.0cm,yshift=-2.75cm]current page.north west)
    {\sffamily\bfseries\fontsize{9}{11}\selectfont SERIE GUÍAS · TECNOLOGÍA E INFORMÁTICA};

  \begin{scope}[shift={([xshift=-3.05cm,yshift=-2.55cm]current page.north east)}]
    \fill[white,opacity=.80] (-.62,-.52) rectangle (.62,.52);
    \draw[milcGradeText,line width=1pt] (-.62,-.52) rectangle (.62,.52);
    \fill[milcGradeDark] (-.42,-.38) rectangle (-.22,.06);
    \fill[milcGradeAccent] (-.10,-.38) rectangle (.10,.24);
    \fill[milcGradePrimary!60!black] (.22,-.38) rectangle (.42,.38);
  \end{scope}

  \node[anchor=west,text=milcGradeText] at ([xshift=1.9cm,yshift=-12.85cm]current page.north west)
    {\sffamily\bfseries\fontsize{124}{124}\selectfont 8};
  \draw[milcGradeText,line width=1.7pt]
    ([xshift=4.52cm,yshift=-11.68cm]current page.north west) circle (.13cm);

  \node[anchor=west,text=milcGradeText,text width=8.7cm,align=left] at ([xshift=10.35cm,yshift=-11.6cm]current page.north west)
    {
     {\sffamily\bfseries\fontsize{19}{22}\selectfont Seguridad en la red II: sexting y grooming — estrategias de protección y conciencia digital}\\[4mm]
     {\sffamily\bfseries\fontsize{23}{26}\selectfont Guía 27-8-TIC}\\[2mm]
     {\sffamily\fontsize{13}{16}\selectfont Periodo Tercero}\\[5mm]
     {\sffamily\bfseries\fontsize{11}{13}\selectfont Institución Educativa Sor María Juliana}\\[1mm]
     {\sffamily\fontsize{10}{12}\selectfont Autor: Dr. Álvaro Cárdenas Orozco}
    };

  \path[fill=milcGradeSoft,rounded corners=7mm]
    ([xshift=1.35cm,yshift=.85cm]current page.south west) rectangle
    ([xshift=-1.35cm,yshift=3.15cm]current page.south east);
  \draw[milcGradeDark,line width=.65pt,rounded corners=7mm]
    ([xshift=1.35cm,yshift=.85cm]current page.south west) rectangle
    ([xshift=-1.35cm,yshift=3.15cm]current page.south east);
  \node[anchor=center,text=milcNegro,text width=17.0cm,align=center] at ([yshift=2.0cm]current page.south)
    {\sffamily\fontsize{10}{12}\selectfont Metodología MILC · Apertura ancestral · Pensamiento computacional · Triángulo Dussel-Estoicismo-Floridi\\
    \textcolor{milcGradeDark}{\bfseries Producto:} Mapa visual "qué hacer ante grooming" + 3 mensajes plantilla para responder/cortar la comunicación + lista de líneas de ayuda.};
\end{tikzpicture}
\mbox{}
\newpage

\section*{Apertura: Seguridad en la red II: sexting y grooming — estrategias de protección y conciencia digital}

\begin{softbox}{milcGradeDark}{milcGradeSoft}{Saber ancestral}
La \textbf{advertencia de los mayores sobre el extraño en la plaza}: \textit{``no aceptes regalos de quien no conoces''}, \textit{``si alguien te ofrece dulces y quiere llevarte aparte, corre y avisa''}. La sabiduría comunitaria protegía a los niños del peligro acechante. Hoy el extraño llega por DM, juego en línea, comentario en TikTok. La sabiduría es la misma; el formato cambió.
\end{softbox}

\begin{softbox}{milcTurquesa}{milcGris}{Saber contemporáneo}
El \textbf{grooming} es la manipulación gradual de un adulto que se hace pasar por igual para abusar de menores. El \textbf{sexting} es el envío de contenido íntimo. Ambos pueden combinarse en patrones de abuso. La protección requiere reconocer banderas rojas, no enviar nunca contenido íntimo, y hablar con adulto de confianza.
\end{softbox}

\begin{softbox}{milcMostaza}{milcGris}{Pregunta-puente}
\textit{¿Qué de la advertencia ancestral —no aceptes regalos del extraño— podemos llevar a las redes y juegos en línea?}
\end{softbox}

\begin{softbox}{milcVerde}{milcGris}{Saber y saber hacer}
Al terminar podrás \textbf{reconocer} grooming en sus etapas tempranas, \textbf{evitar} prácticas de sexting riesgoso, \textbf{proteger} tu privacidad digital, \textbf{actuar} ante intentos de abuso, y \textbf{conocer} la Línea 141 del ICBF para casos urgentes.
\end{softbox}

\small
\begin{tabularx}{\textwidth}{>{\bfseries}p{3.0cm}Y}
\rowcolor{milcGradePrimary!12}Autor & Dr. Álvaro Cárdenas Orozco\\
DBA & Reconoce y previene grooming/sexting aplicando rutas de protección y autocuidado digital.\\
Referentes & MEN Tecnología · ICBF Línea 141 · Modelo MILC · Convención Derechos del Niño · Filosofía de la Liberación (Dussel) · Estoicismo · Infoética (Floridi) · Ley 1336 de 2009 (delitos sexuales online)\\
Producto final & Mapa visual "qué hacer ante grooming" + 3 mensajes plantilla para responder/cortar la comunicación + lista de líneas de ayuda.\\
\end{tabularx}
\normalsize

\begin{titlebox}{milcGradeDark}
Ruta MILC: así vamos a aprender
\end{titlebox}
\begin{tabularx}{\textwidth}{>{\centering\arraybackslash}X>{\centering\arraybackslash}X>{\centering\arraybackslash}X>{\centering\arraybackslash}X}
\phasecard{1}{milcVerde}{milcGris}{Escuta\\[-1mm]\small Observo, escucho y pregunto.} &
\phasecard{2}{milcTurquesa}{milcGris}{Sistematización\\[-1mm]\small Pensamiento computacional.} &
\phasecard{3}{milcMagenta}{milcGris}{Praxis\\[-1mm]\small Construyo y aplico.} &
\phasecard{4}{milcVino}{milcGris}{Evaluación\\[-1mm]\small Reflexión liberadora.}
\end{tabularx}


\begin{titlebox}{milcVerde}
Fase 1 · Escuta
\end{titlebox}

\begin{softbox}{milcVerde}{milcGris}{Escena para escuchar}
Camila recibe mensajes de un \textit{``amigo de un amigo''} en Instagram. Al inicio él la trata bien, le dice cumplidos, le manda música. Después de 2 semanas le pide fotos a cambio de regalos virtuales (skins de juegos). Le dice que es \textit{``secreto entre los dos''}. Camila siente que algo está raro pero no sabe qué. Esos son los pasos clásicos del grooming. ¿Cómo identificarlos a tiempo?
\end{softbox}

\checkbox Reconozco situaciones donde un adulto se hizo pasar por un menor en redes (mías o de conocidos).\\
\checkbox Identifico la diferencia entre amistad genuina y manipulación gradual.\\
\checkbox Distingo entre privacidad de cuenta básica y privacidad reforzada.

\textbf{Una vez (yo o alguien cercano) recibimos mensajes raros de un extraño en redes. Lo que pasó fue:}\\[1mm]
\linead{\linewidth}\\[1mm]
\linead{\linewidth}

\textbf{Lo que distinguió esa situación de una amistad genuina fue:}\\[1mm]
\linead{\linewidth}\\[1mm]
\linead{\linewidth}

\textbf{La principal lección que saqué fue:}\\[1mm]
\linead{\linewidth}\\[1mm]
\linead{\linewidth}

\begin{infoband}{milcVerde}
\textbf{Saber más.} El grooming es DELITO en Colombia (Ley 1336 de 2009). Las penas son altas. Pero la mejor protección es PREVENCIÓN: privacidad de cuentas, no aceptar desconocidos, no enviar fotos íntimas a NADIE. Una foto íntima en internet puede no desaparecer nunca.
\end{infoband}

\begin{titlebox}{milcTurquesa}
Fase 2 · Sistematización con pensamiento computacional
\end{titlebox}

\begin{softbox}{milcTurquesa}{milcGris}{Cuatro pilares aplicados al tema}
El pensamiento computacional descompone el grooming en sus etapas + banderas rojas + protocolos de protección. Reconoce patrones, abstrae la regla \textit{anticipar el patrón salva} y diseña un algoritmo de defensa.
\end{softbox}

\small
\begin{tabularx}{\textwidth}{>{\bfseries\RaggedRight\arraybackslash}p{3.6cm}Y}
\rowcolor{milcTurquesa!90}\color{white}Pilar & \color{white}\bfseries Aplicación al tema\\
1. Descomposición & Descomponer las etapas del grooming: 1) contacto inicial · 2) construcción de confianza · 3) aislamiento ("esto es secreto") · 4) petición íntima · 5) chantaje.\\
2. Reconocimiento de patrones & Banderas rojas: pide foto, secreto, encuentro presencial, regalos virtuales/físicos a cambio de algo, halagos excesivos al inicio.\\
3. Abstracción & Abstracción: si UN paso te incomoda, hay grooming. NO esperes a la etapa final — la prevención es en las primeras señales.\\
4. Algoritmo conceptual & Algoritmo de defensa: \textbf{1)} privacidad cerrada · \textbf{2)} no aceptar desconocidos · \textbf{3)} reconocer bandera roja · \textbf{4)} cortar comunicación · \textbf{5)} hablar con adulto · \textbf{6)} Línea 141 si hay amenaza.\\
\end{tabularx}
\normalsize

\begin{drawbox}{milcTurquesa}{Anatomía del grooming en 5 etapas}
\textbf{1. Contacto:} adulto se hace pasar por igual o conocido. \\
\textbf{2. Confianza:} cumplidos excesivos, regalos virtuales, escucha mucho. \\
\textbf{3. Aislamiento:} \textit{``no le digas a tus papás, esto es secreto''}, \textit{``solo tú me entiendes''}. \\
\textbf{4. Petición íntima:} fotos, videos, encuentros presenciales. A cambio de regalos. \\
\textbf{5. Chantaje:} \textit{``si no me mandas más, mando lo anterior a tu colegio''}. \\
\textbf{Defensa:} cortar en etapa 1 o 2 (lo más temprano posible).
\end{drawbox}

\begin{softbox}{milcMostaza}{milcGris}{Errores comunes y su corrección}
\begin{itemize}
\item Aceptar a desconocidos por curiosidad $\rightarrow$ privacidad cerrada y SOLO conocidos verificados.
\item Creer en "amigos virtuales" como amistad real $\rightarrow$ amistad presencial primero, virtual segundo.
\item Enviar fotos íntimas (sexting) "de confianza" $\rightarrow$ NUNCA. Una vez enviadas, no se controlan.
\item Guardar el secreto pedido por el extraño $\rightarrow$ secreto = bandera roja MÁXIMA. Cuenta a adulto SIEMPRE.
\item No reportar "para no causar problemas" $\rightarrow$ reportar PROTEGE a otros menores que pueden caer.
\end{itemize}
\end{softbox}

\textbf{La bandera roja más temprana que reconocería sería:}\\[1mm]
\linead{\linewidth}\\[1mm]
\linead{\linewidth}

\textbf{Mi mensaje de respuesta (sin escalada) sería:}\\[1mm]
\linead{\linewidth}\\[1mm]
\linead{\linewidth}

\begin{titlebox}{milcMagenta}
Fase 3 · Praxis con pensamiento computacional
\end{titlebox}

\begin{softbox}{milcMagenta}{milcGris}{Construyendo el producto con los cuatro pilares}
Construir un mapa visual de defensa contra grooming es ejercicio de protección comunitaria. Decomponemos las etapas, reconocemos patrones, abstraemos las banderas rojas y producimos material que puede salvar a otros.
\end{softbox}

\small
\begin{tabularx}{\textwidth}{>{\bfseries\RaggedRight\arraybackslash}p{3.6cm}Y}
\rowcolor{milcMagenta!90}\color{white}Pilar & \color{white}\bfseries Aplicación al producto\\
Descomposición & Mapa: 1) etapas del grooming · 2) banderas rojas · 3) qué hacer en cada etapa · 4) mensajes plantilla de corte · 5) líneas de ayuda.\\
Patrones & Patrones de mapas eficaces: visual, simple, números siempre presentes, sin estigmatizar a víctima.\\
Abstracción & Función esencial: que un compañero al ver el mapa identifique si está en grooming y sepa cómo cortar.\\
Algoritmo de armado & Pasos: investigar Ley 1336 + recursos ICBF $\rightarrow$ identificar etapas $\rightarrow$ diseñar visual $\rightarrow$ probar con un compañero $\rightarrow$ revisar con orientador.\\
\end{tabularx}
\normalsize

\begin{drawbox}{milcMagenta}{Checklist del mapa de defensa}
\checkbox 5 etapas del grooming claramente identificadas.\\
\checkbox Banderas rojas en cada etapa.\\
\checkbox 3 mensajes plantilla para cortar comunicación.\\
\checkbox Línea 141 ICBF visible.\\
\checkbox Política de privacidad de Instagram/TikTok citada.\\
\checkbox Diseño visual profesional sin alarmismo.
\end{drawbox}

\begin{softbox}{milcMagenta}{milcGris}{Plantilla de trabajo}
\textbf{Mapa de defensa:} 5 etapas del grooming en infografía visual. \\
\textbf{Etapa 1 (contacto):} \textit{``Si un desconocido te escribe, NO respondas. Bloquea.''} \\
\textbf{Etapa 2 (confianza):} \textit{``Si te llena de halagos al inicio, ya es bandera roja''.} \\
\textbf{Mensaje de corte:} \textit{``No quiero seguir esta conversación. Voy a contarle a mi familia''.} \\
\textbf{Recurso clave:} Línea 141 ICBF gratuita 24/7.
\end{softbox}

\textbf{Mi mapa visual de defensa contra grooming:}\\[1mm]
\linead{\linewidth}\\[1mm]
\linead{\linewidth}\\[1mm]
\linead{\linewidth}

\begin{titlebox}{milcMostaza}
Producto de la sesión
\end{titlebox}
\textbf{Mapa visual + 3 mensajes plantilla + líneas de ayuda}.
\begin{softbox}{milcMostaza}{milcGris}{Criterios de calidad}
\begin{itemize}
\item 5 etapas del grooming.
\item Banderas rojas claras.
\item 3 mensajes plantilla.
\item Líneas de ayuda visibles.
\item Diseño profesional sin alarmismo.
\end{itemize}
\end{softbox}

\begin{titlebox}{milcVino}
Fase 4 · Evaluación liberadora — cinco dimensiones
\end{titlebox}

\begin{softbox}{milcVino}{milcGris}{Sentido de la evaluación}
Evaluar no es solo revisar una nota. Es reconocer qué aprendí, qué cambió en mí, qué emociones manejé, cómo me relaciono con la ciudad y la comunidad, y qué le dejaré a quien viene detrás.
\end{softbox}

\textbf{Mi reflexión liberadora en cinco dimensiones.} Completa cada frase iniciada:

\small
\begin{tabularx}{\textwidth}{>{\bfseries\RaggedRight\arraybackslash}p{3.4cm}Y>{\RaggedRight\arraybackslash}p{4.6cm}}
\rowcolor{milcVino}\color{white}Dimensión & \color{white}\bfseries Mi respuesta (frame starter) & \color{white}\bfseries Algo que descubrí\\
1. Personal & Lo que más cambió en mí fue \linead{6.5cm} & \linead{4.2cm}\\
2. Emocional & La emoción más fuerte fue \linead{2.5cm} y la regulé \linead{3cm} & \linead{4.2cm}\\
3. Ciudadana & En democracia esto importa porque \linead{6cm} & \linead{4.2cm}\\
4. Local (Cartago) & En mi barrio sería útil para \linead{6.5cm} & \linead{4.2cm}\\
5. Intergeneracional & A mi abuela/abuelo le diría \linead{6.5cm} & \linead{4.2cm}\\
\end{tabularx}
\normalsize

\begin{infoband}{milcVino}
\textbf{Para la próxima sustentación.} Vuelve a esta tabla en seis meses. Verás que las respuestas cambian — no porque lo aprendido cambie, sino porque tú cambias. La evaluación liberadora es una conversación contigo a través del tiempo.
\end{infoband}


\begin{titlebox}{milcGradeDark}
Triángulo de pensamiento · cierre filosófico
\end{titlebox}

\begin{softbox}{milcVino}{milcGris}{Dussel · lente del nosotros}
\textit{``Sin reconocimiento del Otro no hay justicia, sólo poder.''}

El groomer ejerce poder sobre el menor borrando su voz. La justicia digital empieza por reconocer al Otro (al menor) como sujeto con dignidad propia. La Ley 1336 colombiana es esa justicia legal.

\textbf{¿Tu mapa restituye dignidad a la víctima o la convierte en "caso difícil"?}
\end{softbox}
\linead{\linewidth}\\[1mm]
\linead{\linewidth}

\begin{softbox}{milcMostaza}{milcGris}{Estoicismo · Marco Aurelio · cuidado interior}
\textit{``Sé como el promontorio donde rompen las olas, firme y tranquilo.''}

El groomer busca quebrar la firmeza interior del menor con halagos y manipulación. La firmeza estoica del menor es saber que su valor NO depende de la opinión del agresor.

\textbf{¿Qué firmeza interior puedes cultivar para que ningún halago externo te haga ceder?}
\end{softbox}
\linead{\linewidth}\\[1mm]
\linead{\linewidth}

\begin{softbox}{milcTurquesa}{milcGris}{Floridi · lente de la infoesfera}
\textit{``El cuidado de los datos es una forma contemporánea de cuidar a las personas.''}

La privacidad digital del menor (cuenta cerrada, datos protegidos) es protección legal y emocional. Cuidar el archivo digital del joven es cuidar al joven mismo.

\textbf{¿Tu mapa enseña no solo a reconocer grooming sino a PROTEGER datos antes del riesgo?}
\end{softbox}
\linead{\linewidth}\\[1mm]
\linead{\linewidth}

\begin{titlebox}{milcVino}
Compromiso final
\end{titlebox}

\small
\begin{tabularx}{\textwidth}{>{\RaggedRight\arraybackslash}p{3.0cm}YYY}
\rowcolor{milcVino}\color{white}\bfseries Criterio & \color{white}\bfseries Lo logré & \color{white}\bfseries En proceso & \color{white}\bfseries Necesito apoyo\\
Comprensión del tema & Explico ideas centrales con mis palabras. & Reconozco ideas pero falta precisión. & Necesito repasar conceptos base.\\
Pensamiento computacional & Apliqué los 4 pilares al diseñar mi producto. & Apliqué algunos pilares. & Necesito releer la fase 2.\\
Aplicación y producto & Mi producto cumple los criterios y se puede revisar. & Producto funcional con ajustes pendientes. & Aún en construcción.\\
Reflexión 5 dimensiones & Respondí cada dimensión con honestidad. & Respondí algunas con detalle. & Mis respuestas son muy generales.\\
Triángulo de pensamiento & Conecté las 3 voces con mi trabajo real. & Conecté al menos una. & No relaciono las voces con mi proceso.\\
\end{tabularx}
\normalsize

\textbf{Hoy entendí que:}\\[1mm]
\linead{\linewidth}\\[1mm]
\linead{\linewidth}

\textbf{Mi compromiso para la próxima sesión es:}\\[1mm]
\linead{\linewidth}\\[1mm]
\linead{\linewidth}

\begin{softbox}{milcMostaza}{milcGris}{Cita de despedida · Marco Aurelio}
\textit{``El obstáculo en el camino se vuelve el camino. Lo que estorba al camino se vuelve camino.''}\\[1mm]
Cada error de hoy, bien observado, se convierte en pieza del camino que pisarás mañana.
\end{softbox}

\small
\begin{tabularx}{\textwidth}{>{\bfseries}p{3.6cm}Y}
\rowcolor{milcGradeDark!12}Nombre completo: & \linead{10cm}\\
Fecha: & \linead{4cm} \quad Curso/grupo: \linead{4cm}\\
Mi firma: & \linead{9cm}\\
\end{tabularx}
\normalsize

\begin{infoband}{milcGradeDark}
\textbf{Próxima sesión.} Aplicaremos la estética de la liberación a contenidos digitales: belleza de la cultura popular como acto de resistencia y dignidad.
\end{infoband}

\end{document}
